\documentclass[11pt]{article}
\usepackage{icmlko}
\begin{document}

\title{모형을 고치면 같이 가고 자료를 고치면 갈라선다 --- 한 노트 뒤에 그 노트가 반쯤 뒤집혔다}
\author{월드모델 실험실}
\date{노트 570 · 2026년 8월 5일}
\maketitle

\begin{abstract}
노트 $569$ 가 이 세션의 팔 $15$ 개를 모아 ``\textbf{판은 전이를 예측하지
못한다}''($\rho \approx +0.1\!\sim\!0.3$)를 냈다. 그리고 \textbf{한계}를
적어 뒀다 --- ``팔 $15$ 개가 전부 자료 계열이다''. 원장에서 옛 팔을 긁어
\textbf{모형 손잡이 계열}을 채우니 부호가 갈렸다.

\begin{center}
\begin{tabular}{lrr}
\hline
계열 & $n$ & 판 $\leftrightarrow$ 전이 순위상관 \\
\hline
\textbf{모형 손잡이} & $17$ & $\mathbf{+0.591}$ \\
\textbf{자료} & $15$ & $\mathbf{-0.089}$ \\
\hline
\end{tabular}
\end{center}

극단값(노트 $544$ 의 재앙)을 빼면 자료 계열이 $\mathbf{-0.341}$ 로 더
선명해지고, 모형 계열은 $+0.579$ 로 그대로다. \textbf{부호가 반대다.}
\end{abstract}

\begin{figure}[t]
\centering
\includegraphics[width=\textwidth]{figs/f1.pdf}
\caption{\textbf{가} --- 모형 손잡이. \textbf{나} --- 자료.
\textbf{다} --- 같은 관문, 다른 세계.}
\end{figure}

\section{읽기}

\begin{itemize}
\item \textbf{모형을 더 잘 만들면 집 안과 집 밖이 같이 좋아진다} ---
모형은 \emph{공유된 것}을 배운다. 깊이 $\cdot$ 학습률 $\cdot$ 반복 $\cdot$
행 무게를 고친 팔 $17$ 개가 그렇다(노트 $456\!\sim\!518$).
\item \textbf{자료를 바꾸면 하나를 올릴 때 다른 하나가 내려가는 경향}이다
--- 자료 변경은 \emph{도메인 특유의 것}을 넣거나 뺀다. 축 가림 $\cdot$
\texttt{grp} $\cdot$ category $\cdot$ 검색 $\cdot$ 달력 팔 $15$ 개가 그렇다.
\end{itemize}

\section{노트 505 와 맞물린다}

노트 $505$ 가 ``\textbf{모형 갈래는 닫혔다} --- 남은 여지 $+0.0009$ 대 자료
손잡이 $0.0283$''을 적었다. 둘을 겹치면

\begin{center}
\begin{tabular}{lll}
\hline
 & 판과 전이 & 남은 여지 \\
\hline
모형 손잡이 & \textbf{같이 간다}($+0.58$) & \textbf{없다}($+0.0009$) \\
자료 & \textbf{갈라선다}($-0.31$) & \textbf{크다}($0.0283$) \\
\hline
\end{tabular}
\end{center}

\begin{quote}
\textbf{쉬운 쪽은 이미 끝났고, 남은 쪽은 판이 목표를 못 맞히는 곳이다.}
그것이 이 판이 지금 서 있는 자리다.
\end{quote}

\section{실무 규칙}

\textbf{자료 변경은 판만 보고 고르면 안 되고 다섯 자를 다 봐야 한다.}
모형 손잡이는 판만 봐도 대체로 맞는다. 지금 관문이 둘 다에 다섯 자를
요구하는데, \textbf{자료 쪽에서 그것이 필수}다.

이 사이클의 채택 둘이 그 증거다 --- 노트 $553$(만화 가림: 판 $-0.0001$,
KR $+0.0082$, 앱 $+0.0022$)과 노트 $564$(게임 zero: 판 $+0.0005$, 게임
$40/40$). \textbf{둘 다 판이 아니라 다른 자로 채택됐다.}

\section*{한계를 적어 둔 것이 값을 했다}

노트 $569$ 의 결론은 \textbf{자료 계열에 한정}된다. 그때 ``전부 자료
계열이다''라고 적어 두지 않았다면 ``판은 전이를 예측 못 한다''가 전체
법칙으로 남았을 것이다.

\begin{quote}
\textbf{한계는 다음 노트의 사전 등록이다.} 노트 $569$ 의 남은 것 첫 줄이
``모형 손잡이 계열로도 재 본다''였고, 그것이 하루도 안 지나 그 노트를
반쯤 뒤집었다.
\end{quote}

\section*{조항}

\begin{quote}
\textbf{``판이 전이를 맞히나''는 계열마다 다르게 답해야 한다.} 모형 쪽에서
$+0.58$, 자료 쪽에서 $-0.31$ 이다. \textbf{한 숫자로 답하면 둘 다 틀린다.}
\end{quote}

\section*{남은 것}

\begin{itemize}
\item \textbf{자료 계열의 $-0.31$ 이 왜 음수인가} --- ``도메인 특유의 것을
넣으면 집 안이 오르고 집 밖이 내린다''는 읽기는 그럴듯하지만 안 쟀다.
노트 $540$(도메인 지시자) $\cdot$ 노트 $419\!\sim\!423$(표지는 하나만)이
같은 자리를 가리킨다.
\item \textbf{긁는 방식의 한계} --- 원장 기록을 정규식으로 모았으므로 표기가
다른 팔은 빠졌다. 계열 분류도 노트 번호로 한 거친 것이다.
\item \textbf{셋째 집 밖 짝} --- 지금 KR 만화 $\cdot$ 비게임 앱 둘뿐이다.
모형 계열에서 둘의 상관이 $+0.789$, 자료 계열에서 $+0.574$ 로 \emph{자료
쪽에서 더 갈리므로}, 셋째 짝이 있으면 자료 판정이 단단해진다.
\end{itemize}

\end{document}
